% Paper fragment for NAV.cls. Requires amsmath and booktabs; NAV.cls supplies natbib.
\section{Pl\'eiades imagery: morphology and conditional source-location analysis}
\label{sec:pleiades-conditional}

On 23 March 2014 near 04:00 UTC, Pl\'eiades 1A acquired four 0.5~m
southern Indian Ocean scenes. Geoscience Australia analysed imagery supplied
by French military intelligence and CNES at the ATSB's
request~\citep{minchin2017pleiades}. Twelve of at least 70 catalogued objects
were rated 5 (probably not natural objects or features). This did not identify
aircraft debris. Targets 06, 18, 19, 26 and 27 were retained by mask-stability
screening; other rating-5 objects had weaker geometry. Correspondence-qualified
perturbations (baseline mask $\operatorname{IoU}\geq0.25$) set size intervals:
15/18 for object 06 and 18/18 for every other target; three rejects had switched
components. Selection measured segmentation stability, not aircraft
resemblance.

Eight equal-sized families covered flooding-model and random-pose wings; other
777 parts; matched random geometry; ocean controls; the whole fuselage; and
documented-boundary or random-cut sections. Boeing geometry was
used~\citep{boeing2024acap}. Controls represented containers, nets/ropes,
clustered cargo/plastic/timber, wreckage/hulls and fish-aggregating
devices~\citep{hapag2016container,iotc2023observer,lebreton2018plastic}.
\citet{murphy2013structures} gives six 777 section intervals, including BS~2150
at the aft pressure bulkhead. These reproducible boundaries are not asserted
weak points. Reported Asiana separation, local 777 puncture and A320 aft
damage~\citep{aaib2010ba38,ntsb2010hudson,ntsb2014asiana} do not define an
MH370 breakup prior.

Each family contained 3,372 masks: 26,976 masks and 134,880 family--object
comparisons. Fuselage orientation was isotropic; one fifth were fully visible,
otherwise a linear waterline targeted 15--100\% visibility. These were coverage
sensitivities, not probabilities. Eligibility required visible area, oriented
length and width each between the target's lower and upper limits.
Principal-axis-normalised masks were compared under identity and three
reflections. For candidate $c$, target $o$, family $g$, masks $T_o,P_o$,
reflection set $\mathcal F$, and eligible set $E_o$,
\begin{equation}
\begin{aligned}
  \operatorname{IoU}(A,B)
    &=\frac{|A\cap B|}{|A\cup B|},\\
  \ell(c,o)
    &=1-\frac{
      \max_{f\in\mathcal F}\operatorname{IoU}(T_o,f(c))+
      \max_{f\in\mathcal F}\operatorname{IoU}(P_o,f(c))
    }{2},\\
  \ell^{*}_{g,o}
    &=\min_{c\in g\cap E_o}\ell(c,o).
\end{aligned}
\label{eq:pleiades-mask-loss}
\end{equation}
Lower loss is greater overlap, not an identity probability.

Extended minima were an ocean-control fishing-net-or-rope mask (06), matched
random geometry (18/19), a random-pose wing (26), and documented section 47
(27). Object-06 ocean and random-geometry losses differed by 0.0007. Section 47
scored 0.3058 versus 0.3145 for random geometry, but 0.3556 with fully visible
masks. The whole shell ranked lowest for none; its minima required axis/view
cosines 0.885--0.988 and visible fractions 0.305--0.672: near-end-on, occluded,
or both. Random geometry matched dimension ranges, not outlines. 777 shapes
were admissible but not preferentially selected; morphology remained
non-specific.

\begin{table}[t]
  \centering
  \footnotesize
  \caption{Equal-budget Pl\'eiades morphology controls. Lower loss is better.}
  \label{tab:pleiades-morphology-controls}
  \begin{tabular}{lrrr}
    \toprule
    Family & Eligible objects & Mean lowest loss\textsuperscript{a} &
      Objects at lowest loss \\
    \midrule
    Flooding-model wing            & 1/5 & 0.3286 & 0 \\
    Random-pose wing               & 5/5 & 0.3232 & 1 \\
    Other shortlisted 777 parts    & 4/5 & 0.3704 & 0 \\
    Matched random geometry        & 5/5 & 0.3267 & 2 \\
    Generic ocean objects          & 4/5 & 0.3532 & 1 \\
    Whole 777 fuselage               & 4/5 & 0.3472 & 0 \\
    Documented 777 fuselage sections & 5/5 & 0.3318 & 1 \\
    Random 777 fuselage sections     & 4/5 & 0.3388 & 0 \\
    \bottomrule
  \end{tabular}

  \vspace{2pt}
  \parbox{0.94\linewidth}{\footnotesize\textsuperscript{a}Mean over eligible
  objects only; family means are therefore not directly comparable.}
\end{table}

For conditional source location, \citet{griffin2017drift} propagated debris
forward from trial sites assuming some objects were from 9M-MRO. A separate
clean-sheet forward-ensemble reconstruction released particles from 5,289
cells along and within $\pm100$~NM of the seventh arc. BRAN2016 and OSCAR v2
were evaluated separately~\citep{dohan2021oscar}, with common NCEP
winds~\citep{kalnay1996reanalysis}, three windage alternatives and 5~NM/day
dispersion. All twelve location likelihoods were averaged equally; no object
was randomly drawn. The five-object set was a sensitivity.

BRAN's mode was 35.418$^{\circ}$S, 92.886$^{\circ}$E with a 24,309~km$^2$ 90\%
region; OSCAR's was 34.917$^{\circ}$S, 92.047$^{\circ}$E with 18,393~km$^2$.
Averaging uniform subsets of one to five objects reproduces that PDF for any
subset-size prior; a product needs correlated debris and detection models. The
50:50 BRAN--OSCAR prior mixture has mode 35.300$^{\circ}$S,
92.891$^{\circ}$E and a 25,896~km$^2$ 90\% region. It is not
independent-evidence multiplication or a calibrated posterior. These are
conditional hypotheses for an integrated estimator; none establishes identity
or updates the main estimate unconditionally.
